\chapter
[\CO{[GEF]} Gefahrenzonen, Gefahrentransporte, Gefahrengutunfall]
{\CC{[GEF]}\\Gefahrenzonen, \\Gefahrentransporte,
\\Gefahrengutunfall}
\label{topic:gefahrengutunfall}

\ChapterIntro
{%
~~~
}
{%
\begin{description}
  \item[Maintainer:] Roman Koch
  \item[Autoren:] Diverse
  \item[Reviewer:] Gerald Höritzmiller
  \item[Version:] {\DocVersion} ({\DocVersionInfo})
  \item[SHA1:] {\DocGitCommit}
\end{description}

{\TxTeilkapitelAASS}
}

% TODO : Layout-Anpassung
\TODO{GABS}{yyyy-mm-dd}{Layout-Anpassung notwendig.}{Dieser Teil entspricht nicht dem Standard-Layout!}


% \noindent\Parallel{
% % \includegraphics[width=\linewidth]{./images/noauth/AASdev20090228-img149.jpg}
% \includegraphics[width=\linewidth]{./icons/under-construction.png}
% }{
% % \includegraphics[width=\linewidth]{./images/noauth/AASdev20090228-img150.jpg}
% \includegraphics[width=\linewidth]{./icons/under-construction.png}
% }

\TODO{GABS}{2010-04-27}{Bilder 2x Gefahrengutunfall/LKW}{}

\section{Gefahrenzonen}
\label{topic:gefahrenzonen}

\Layout{0}{Mögliche Gefahren für den Sanitäter}{
\begin{itemize}
    \item Brand bzw. Explosionsgefahr
    \item Entweichen unter Druck stehender Gase
    \item Giftigkeit, Ansteckungsgefahr
    \item Ätzwirkung
    \item Radioaktivität
    \item Umweltgefährdung
\end{itemize}
}{
\begin{itemize}
    \item Brand bzw. Explosionsgefahr
    \item Entweichen unter Druck stehender Gase
    \item Giftigkeit, Ansteckungsgefahr
    \item Ätzwirkung
    \item Radioaktivität
    \item Umweltgefährdung
\end{itemize}
}{}{}{}




\Layout{0}{Gefahrenquellen können entstehen bei}{
\begin{itemize}
    \item Unfällen bei der Herstellung
    \item Transportunfällen
    \item Unfällen bei der Verwendung von gefährlichen Stoffen
\end{itemize}
}{
\begin{itemize}
    \item Unfälle bei der Herstellung
    \item Transportunfälle
    \item Unfälle bei der Verwendung von gefährlichen Stoffen   
\end{itemize}
}{}{}{}


\Layout{47}{}{}{}{
\noindent\begin{tabulary}{\linewidth}{LL}\toprule
\multicolumn{2}{c}{\TabHeading{Andere Gefahrenquellen}}
\\\midrule
Landwirtschaft
 &
\ce{CO2} in Gärkellern, Siloanlagen, Brunnenschächte

Pflanzenschutzmittel
\\\addlinespace
Klärgruben
 &
\ce{CO2}, Methan, Schwefelwasserstoff
\\\addlinespace
Obstlagerhallen
 &
wenig \ce{O2}, viel \ce{CO2} oder \ce{N2}, halten Obst frisch
\\\addlinespace
Brände in Wirtschaftsgebäuden
 &
extrem gefährliche Rauchgase
\\\addlinespace
Tiefgaragen und Tunnels
 &
erhöhte \ce{CO}- und \ce{CO2}-Werte
\\\addlinespace
Flüssiggas
 &
Campingplätze, feste und mobile Küchen
\\\addlinespace
Pools, Schwimmbäder
 &
Chlor (Wasserdesinfektion)
\\\addlinespace
Industrie
 &
Kühlmittel in großen Kälteanlagen (Ammoniak)
\\\addlinespace
Speziallabors, Einrichtungen mit und zur Verarbeitung von infektiösen
Abfällen oder Tierkadaver
 &
Infektionsgefahr
\\\addlinespace
Nuklearmedizinische Institute, Industrie, Forschungseinrichtungen
 &
Radioaktivität
\\\bottomrule
\end{tabulary}
}{}{}



\Layout{2}{Strom -- allgemein}{%
\index{Gefahrenbereich!Strom}\label{topic:gefahrenbereich-strom}
%
Eine Starkstrom- bzw. Hochspannungsanlage muss durch einen \EE{Fachmann}
(E-Werk, ÖBB, Betriebsleitung, ...) abgeschalten bzw. geerdet werden. 

Die Anlage muss
gegen Wiedereinschalten gesichert werden! Patient mit isolierendem
Material aus dem Stromkreis befreien. Ein Sicherheitsabstand zu
stromführenden Leitungen muss eingehalten werden! 
}{%

}{}{}{}

\Layout{2}{Freileitungen}{%
Bei Freileitungen mit
220\,000 Volt (220\,\MetricUnit{kV}) sind mindestens 4\Meter* Abstand zu halten.
\cite{Jaros1} 

Bei herabhängenden Hochspannungsleitungen oder
nach Blitzeinschlägen auf die Schrittspannung achten: Die
Spannung fällt im Boden nicht sofort auf 0~Volt ab, innerhalb
eines gewissen Umkreises kann daher zwischen den beiden
Beinen am Boden eine so genannte Schrittspannung auftreten.
Befinden Sie sich innerhalb eines solchen Gebietes lassen Sie
Ihre Beine eng beisammen und springen aus diesem Bereich
heraus. Je nach Bodenbeschaffenheit können Schrittspannungen
bis zu einer Distanz von 20\Meter* von der herunterhängenden
Hochspannungsleitung auftreten. \cite{RedelsteinerHRNS1} 
}{%

}{}{}{}

\Layout{0}{Bahnanlagen}{%

\begin{description}
  \item[ÖBB-Einsatzleiter] Der \EE{ÖBB-Einsatzleiter} ist
  der Ansprechpartner für die Sicherheit im Gleisbereich!
  Einsätze im Gleisbereich dürfen grundsätzlich nur mit
  Zustimmung des ÖBB-Einsatzleiters erfolgen!
  
  Bei \E{Gefahr in Verzug} dürfen alle erforderlichen
  Maßnahmen zur Menschenrettung nur unter Berücksichtigung
  der Eigensicherung (Einhaltung von Schutzabständen \ldots)
  und der Eigenverantwortung gesetzt werden.
  \cite{OeBBRettEinGleis2009}
  \item[Problem Zugverkehr] Bei
  Notfällen auf Bahngleisen ist auf jeden Fall der
  Streckenabschnitt von der ÖBB sperren zu lassen (über die
  Leitstelle)!
  \item[Problem Lokalisation] An den Masten der
  Oberleitungen der ÖBB sind Schilder mit Kilo- und
  Hektometerangaben angebracht\ZFootnote{Die obere Zahl gibt
  den Kilometer, die untere den Hektometer (100\Meter*)
  an.}. Diese müssen bei der Leitstelle gemeldet werden. Bei
  nicht-elektrifizierten Strecken finden sich die Angaben
  auf Hektometersteinen am Rand der Strecke. Weiters verfügt
  jeder Mast über eine Mastnummer, auch diese kann angegeben
  werden.
  \cite{OeBBRettEinGleis2009}
  
  \Fazit{Die Nummernsysteme der ÖBB sind eine eigene Wissenschaft, im Zweifel sind
  zusätzlich zur Nummer auch anzugeben, woher diese stammt.}
  \item[Problem Strom] Die Eisenbahn arbeitet mit einer Oberleitungsspannung von
  15\,000\,\MetricUnit{V}, auch bei gesenktem Stromabnehmer oder Dieselfahrzeugen kann durch Kondensatorladung eine
  Spannung von bis zu 3\,000\,\MetricUnit{V} anliegen.
  
  \Speech{"`Die Freischaltung und Erdung wird vom
  ÖBB-Einsatzleiter veranlasst. Hilfs- und Einsatzkräfte
  erhalten beim ÖBB-Einsatzleiter vor Einsatzbeginn eine
  verbindliche Auskunft über den Schaltzustand der
  Oberleitung im Einsatzbereich. Alle Anlagenteile --
  ausgenommen jene, die vom ÖBB-Einsatzleiter als
  freigeschaltet und geerdet bekanntgegeben werden -- sind
  als unter Spannung stehend zu betrachten!"'}
  \cite{OeBBRettEinGleis2009}
  
  Bei herabhängeenden Leitungen ist ein Sicherheitsabstand
  von mindestens 15\Meter* einzuhalten!
  
  Ähnliches gilt für Stadt- und U-Bahn-Anlagen. Da die
  Stromversorgung oft nicht durch eine Oberleitung, sondern
  durch andere Vorrichtungen in Bodennähe (Stromschiene,
  etc.) erfolgt, besteht auch bei intakten Anlagen die
  Gefahr eines Stromschlags, sofern die Gleisanlagen
  betreten werden. Als Traktionsspannung werden oft
  Spannungen um die 700\,\MetricUnit{V} eingesetzt.
  \cite{U61989}
  \item[Problem Bremsweg]
  Schienenfahrzeuge haben einen \EE{massiv verlängerten
  Bremsweg}: Ein Güterzug benötigt ca. 1\,\MetricUnit{km},
  ein Personenreisezug bis zu \EE{2,6\,\MetricUnit{km}}, um
  zum Stillstand zu kommen!
  
  \item[Problem Gefahrgut]
  Mit der Bahn transportiertes Gefahrgut wird analog zum
  Straßenverkehr mit Gefahrenzahl und UN-Nummer
  gekennzeichnet (\prettyref{topic:gefahrengutunfall}). Die
  zur Fracht gehörigen Frachtbriefe sind zu beachten.
\end{description}
}{%
\begin{itemize}
  \item ÖBB Einsatzleiter
  \item Zugverkehr
  \item Lokalisation: Hektometerschilder, -steine
  \item Strom:
  \begin{itemize}
    \item Bis zu 15\,000\,\MetricUnit{V}
    \item Erdung und Frei gabe durch ÖBB-Einsatzleiter
  \end{itemize}
  \item Bremsweg: ↑↑↑
  \item Gefahrgut
\end{itemize}


}{}{}{}




\Layout{0}{{Linktip}}{%
\begin{labeling}{mmm}
  \item[\cite{OeBBRettEinGleis2009}] \emph{ÖBB Infrastruktur AG: Handbuch
  Rettungseinsatz im ÖBB-Gleisbereich. online, Oktober 2009. \url{http://www.roteskreuz.at/uploads/media/Handbuch_Einsatz_Gleisbereich.pdf}}
\end{labeling}
}{\url{http://www.roteskreuz.at/uploads/media/Handbuch_Einsatz_Gleisbereich.pdf}
}{}{}{}















\section{Kennzeichnung von gefährlichen Stoffen und
Gefahrenbereichen}
\label{topic:gefahrsymbole}
% TODO : Layout-Anpassung
\TODO{GABS}{yyyy-mm-dd}{Layout-Anpassung notwendig.}{Dieser Teil entspricht nicht dem Standard-Layout!}

\Layout{47}{}{}{}{
\begin{center}
\begin{tabulary}{\textwidth}{LLLL}
\toprule
\multicolumn{2}{c}{\TabHeading{Hinweis auf Gefahr}} 
&
\multicolumn{2}{c}{\TabHeading{Warnzeichen}}
\\\midrule

\includegraphics[width=1.475cm,height=1.452cm]{./images/noauth/gefahr-explosion.png}
 &
Explosionsgefährlich
 &
\includegraphics[width=1.541cm,height=1.426cm]{./images/noauth/warnung-allgemein.png}
 &
Allgemeine Gefahr
\\
\includegraphics[width=1.452cm,height=1.414cm]{./images/noauth/gefahr-giftig.png}
 &
Giftig
 &
\includegraphics[width=1.673cm,height=1.414cm]{./images/noauth/warnung-radioaktiv.png}
 &
Radioaktive Stoffe, ionisierende Strahlen
\\
\includegraphics[width=1.497cm,height=1.611cm]{./images/noauth/gefahr-hochentzuendlich.png}
 &
Hochentzündlich
 &
\includegraphics[width=1.696cm,height=1.412cm]{./images/noauth/warnung-gasflaschen.png}
 &
Gasflaschen
\\
\includegraphics[width=1.475cm,height=1.435cm]{./images/noauth/gefahr-reizend.png}
 &
Reizend
 &
\includegraphics[width=1.673cm,height=1.414cm]{./images/noauth/warnung-laser.png}
 &
Laserstrahlen
\\
\includegraphics[width=1.488cm,height=1.599cm]{./images/noauth/gefahr-aetzend.png}
 &
\"Atzend
 &
 &
\\\hline
\end{tabulary}
\end{center}
}{}{}

\subsection{Transport gefährlicher Güter}

Kennzeichnungspflicht besteht erst dann, wenn die transportierte Menge
eines Stoffes die gesetzlich erlaubte Mindestmenge überschreitet.
Diese Mindestmenge gilt für jeden einzelnen Stoff und nicht für die
Gesamtmenge.

Es können also erhebliche Mengen gefährlicher Güter ohne
Kennzeichnung des Fahrzeuges transportiert werden. 

\subsection{Gefahrzettel}
\label{topic:gefahrzettel}

Bei Stückguttransporten wird jedes Versandstück gekennzeichnet. An
Fahrzeugen, die kein Stückgut transportieren, müssen Gefahrenzettel
an den Längsseiten und an der Rückseite angebracht werden.

\Layout{57}{}{}{}{
\begin{tabular}{m{2.cm}m{3.5cm}m{2.cm}m{3.5cm}}
 Beispiele
 &
 &
 &
\\
\includegraphics[width=1.659cm,height=1.885cm]{./images/noauth/gefahrzettel-feuergefaehrlich-fluessig.png}
 &
Feuergefährlich \par(flüssige Stoffe)
 &
\includegraphics[width=1.856cm,height=2.096cm]{./images/noauth/gefahrzettel-selbstentzuendlich.png}
 &
Selbstentzündlich
\\
\includegraphics[width=1.835cm,height=1.902cm]{./images/noauth/gefahrzettel-feuergefaehrlich-fest.png}
 &
Feuergefährlich \par(feste Stoffe)
 &
\includegraphics[width=1.909cm,height=2.162cm]{./images/noauth/gefahrzettel-enzuendlich-wasser.png}
 &
Entzündbare Gase bei Berührung mit \EE{Wasser}
\\
\end{tabular}
}{}{}

\subsection{Gefahrentafel}
\label{topic:warntafel}

\Layout{57}{}{}{}{
\begin{tabular}{m{3.225cm}m{3.225cm}m{3.225cm}m{3.225cm}}
\multicolumn{2}{m{6.65cm}}{\centering Allgemeine Kennzeichnung
(Sammeltransporte)} &
\multicolumn{2}{m{6.65cm}}{\centering Spezielle
Kennzeichnung} 
\\\midrule
\includegraphics[width=\linewidth]{./images/khd/warntafel-allgemein.png}
 &
 &
\includegraphics[width=\linewidth]{./images/khd/warntafel-un-60-1710.png}
 &
Gefahrnummer\par (\T{Kemler-Nummer})

Stoffnummer\par (\T{UN-Nummer})
\\\bottomrule
\end{tabular}
}{Kennzeichnung durch Gefahrentafel.}{}


\subsection{Gefahrnummer}
\label{topic:gefahrnummer}

\Layout{57}{}{}{}{
\begin{tabular}{cp{11cm}}
\addlinespace
 2 & Entweichen von Gas durch Druck oder chemische Reaktion
\\\addlinespace
 3 & Entzündbarkeit von Flüssigkeiten (Dämpfen) und Gasen
\\\addlinespace
 4 & Entzündbarkeit fester Stoffe
\\\addlinespace
 5 & Oxidierende (brandfördernde) Wirkung
\\\addlinespace
 6 & Giftig
\\\addlinespace
 7 & Radioaktivität
\\\addlinespace
 8 & \"Atzwirkung
\\\addlinespace
 9 & Gefahr einer spontan heftigen Reaktion
\\\addlinespace
 {\sffamily X} & Reagiert gefährlich mit Wasser
\\\addlinespace\bottomrule
\end{tabular}
}{Gefahrennummer.}{}


Die Verdopplung einer Ziffer weist auf die Zunahme der entsprechenden 
Gefahr hin.

\Cave{Ein "`{\sffamily X}"'  vor der Nummer zur Kennzeichnung der Gefahr
bedeutet, dass der Stoff in  gefährlicher Weise \Ca{mit Wasser reagiert}!}

Wenn die Gefahr eines Stoffes ausreichend von einer einzigen Ziffer 
angegeben werden kann, wird dieser Ziffer eine Null angefügt.

Folgende Ziffernkombinationen haben eine besondere Bedeutung:

\Layout{57}{}{}{}{
\begin{tabular}{cp{11cm}}
\addlinespace
22 & tiefgekühltes Gas
\\\addlinespace
X323 & entzündbarer flüssiger Stoff, der mit Wasser gefährlich reagiert, wobei entzündbare Gase entweichen
\\\addlinespace
X333 & selbstentzündliche Flüssigkeit, die mit Wasser gefährlich reagiert
\\\addlinespace
X423 & entzündbarer fester Stoff, der mit Wasser gefährlich reagiert, wobei brennbare Gase entweichen
\\\addlinespace
44 & entzündbarer fester Stoff, der sich bei erhöhter Temperatur in Geschmolzenem Zustand befindet
\\\addlinespace
539 & entzündbares organisches Peroxid
\\\addlinespace
90  & verschiedene gefährlich Stoffe
\\\addlinespace\bottomrule
\end{tabular}
}
{Spezielle Ziffernkombinationen bei der Gefahrennummer.}
{}









\subsection{Neue Kennzeichnung}

\begin{WidePar}
\FormatTable

\begin{tabularx}{\linewidth}{XXXXXXXXX}
\toprule\addlinespace
\multicolumn{9}{c}{\TabSub{Gefahrgutsymbole}}
\\\addlinespace
\multicolumn{9}{c}{\Z{\footnotesize nach Verordnung
(EG) Nr. 1272/2008 gültig ab 1.12.2010 \cite{EG-1272-2008-De}}}
\\\addlinespace\midrule\addlinespace

\includegraphics[width=\linewidth]{./images/GEF/gefahrensymbole/ghs01_explodingbomb.png}
&

\includegraphics[width=\linewidth]{./images/GEF/gefahrensymbole/ghs02_flame.png}
&

\includegraphics[width=\linewidth]{./images/GEF/gefahrensymbole/ghs03_flameovercircle.png}
&

\includegraphics[width=\linewidth]{./images/GEF/gefahrensymbole/ghs04_gascylinder.png}
&

\includegraphics[width=\linewidth]{./images/GEF/gefahrensymbole/ghs05_corrosion.png}
&

\includegraphics[width=\linewidth]{./images/GEF/gefahrensymbole/ghs06_skullandcrossbones.png}
&

\includegraphics[width=\linewidth]{./images/GEF/gefahrensymbole/ghs07_exclamationmark.png}
&

\includegraphics[width=\linewidth]{./images/GEF/gefahrensymbole/ghs08_healthhazard.png}
&

\includegraphics[width=\linewidth]{./images/GEF/gefahrensymbole/ghs09_environment.png}
\\\addlinespace
Gefahr:\newline
Explosions\-gefährlich
&
Gefahr:\newline
Leicht-/Hoch\-entzündlich 
&
Gefahr:\newline
Brandfördernd 
&
Achtung:\newline
Komprimierte Gase 
&
Gefahr:\newline
Ätzend 
&
Gefahr:\newline
Giftig/ Sehr giftig 
&
Achtung:\newline
Gesundheits\-gefährdend 
&
Gefahr:\newline
Gesundheits\-schädlich  
&
Warnung:\newline
Umwelt\-gefährdend
\\\addlinespace
\Z{\tiny GHS01}
&
\Z{\tiny GHS02}
&
\Z{\tiny GHS03}
&
\Z{\tiny GHS04}
&
\Z{\tiny GHS05}
&
\Z{\tiny GHS06}
&
\Z{\tiny GHS07}
&
\Z{\tiny GHS08}
&
\Z{\tiny GHS09}
\\\addlinespace\midrule\addlinespace
\multicolumn{9}{c}{Diese GHS-Symbole werden mit
Gefahrenhinweisen und Sicherheitshinweisen ergänzt!}
\\\addlinespace\toprule\addlinespace
\multicolumn{9}{c}{\TabSub{Alte Zeichen und
Gefahrenbezeichnungen} (Bis Mitte 2017)} 
\\\addlinespace\midrule\addlinespace
\includegraphics[width=\linewidth]{./images/GEF/gefahrensymbole/gefahr5.png}
&
\includegraphics[width=\linewidth]{./images/GEF/gefahrensymbole/gefahr1.png}
&
\includegraphics[width=\linewidth]{./images/GEF/gefahrensymbole/gefahr2.png}
&
&
\includegraphics[width=\linewidth]{./images/GEF/gefahrensymbole/gefahr3.png}
&
\includegraphics[width=\linewidth]{./images/GEF/gefahrensymbole/gefahr6.png}
&

&

&
\includegraphics[width=\linewidth]{./images/GEF/gefahrensymbole/gefahr4.png}
\\\addlinespace
\multicolumn{1}{c}{E} 
&
\multicolumn{1}{c}{F, F+} 
&
\multicolumn{1}{c}{O} 
&
&
\multicolumn{1}{c}{C} 
&
\multicolumn{1}{c}{T, T+} 
&

&

&
\multicolumn{1}{c}{N} 
\\\addlinespace
Explosions\-gefährlich 
&
Leicht-/Hoch\-entzündlich 
&
Brandfördernd 
&
keine Entsprechung 
&
Ätzend 
&
Giftig/ Sehr giftig
&
keine Entsprechung
&
keine Entsprechung 
&
Umwelt\-gefährlich
\\\addlinespace\bottomrule
\end{tabularx}
\end{WidePar}




Fahrzeuge, die \E{ab einer bestimmten Menge} Güter befördern, von
denen eine Gefahr für Lebewesen und Umwelt ausgehen kann (Gefahrgüter),
müssen als solche deutlich gekennzeichnet sein.
Die Kennzeichnung ermöglicht den Helfern bei
einem Unfall die richtigen Maßnahmen
treffen zu können 
% (Klassifizierung nach ADR - Übereinkommen
% über die internationale Beförderung gefährlicher Güter auf
% der Straße).








GEFAHRGUT-KENNZEICHNUNG

Fahrzeuge, die ab einer bestimmten Menge Güter befördern von denen eine
Gefahr für Lebewesen und Umwelt ausgehen kann, müssen als solche
deutlich gekennzeichnet sein.

Die Kennzeichnung erfolgt insbesondere deshalb, damit bei einem Unfall
die Hilfsmannschaften die richtigen Maßnahmen treffen können
(Klassifizierung nach ADR - Übereinkommen über die internationale
Beförderung gefährlicher Güter auf der Straße).

~


\bigskip



\begin{supertabular}{m{0.6\linewidth}m{0.4\linewidth}}
\hline
Gefahrklasse 
&
Symbol

(Auswahl)
\\\hline
Klasse 1

Explosive Stoffe und Gegenstände mit Explosivstoff

Unterklassen:\newline
1.1, 1.2, 1.3, 1.4, 1.5, 1.6 
&
~\newline
\includegraphics[width=1.5cm]{./images/GEF/gefahrgut/gefahrgut_1.png}
\\\addlinespace\midrule\addlinespace
Klasse~2.1

Entzündbare Gase 
&
~\newline
\includegraphics[width=1.5cm]{./images/GEF/gefahrgut/gefahrgut_21.png}
 \\\hline
Klasse~2.2

Nicht entzündbare,\newline
nicht giftige Gase 
&
~\newline
\includegraphics[width=1.5cm]{./images/GEF/gefahrgut/gefahrgut_22.png}
\\\hline
Klasse~2.3

Giftige Gase 
&
~\newline
\includegraphics[width=1.5cm]{./images/GEF/gefahrgut/gefahrgut_23.png}
 \\\hline
Klasse~3

Entzündbare\newline
flüssige Stoffe 
&
~\newline
\includegraphics[width=1.5cm]{./images/GEF/gefahrgut/gefahrgut_3.png}
\\\hline
Klasse~4.1

Entzündbare feste Stoffe, selbstzersetzliche Stoffe und
desensibilisierte explosive feste Stoffe 
&
~\newline
\includegraphics[width=1.5cm]{./images/GEF/gefahrgut/gefahrgut_41.png}
\\\addlinespace\midrule\addlinespace
Klasse~4.2

Selbstentzündliche Stoffe 
&
~\newline
\includegraphics[width=1.5cm]{./images/GEF/gefahrgut/gefahrgut_42.png}
\\\addlinespace\midrule\addlinespace
Klasse~4.3

Stoffe, die in Berührung mit Wasser entzündliche Gase entwickeln 
&
~\newline
\includegraphics[width=1.5cm]{./images/GEF/gefahrgut/gefahrgut_43.png}
\\\addlinespace\midrule\addlinespace
Klasse~5.1

Entzündend\newline
(oxidierend)\newline
wirkende Stoffe 
&
\includegraphics[width=1.5cm]{./images/GEF/gefahrgut/gefahrgut_51.png}
\\\addlinespace\midrule\addlinespace
Klasse~5.2

Organische Peroxide
&
~\newline
\includegraphics[width=1.5cm]{./images/GEF/gefahrgut/gefahrgut_52.png}
\\\addlinespace\midrule\addlinespace
Klasse~6.1

Giftige Stoffe 
&
~\newline
\includegraphics[width=1.5cm]{./images/GEF/gefahrgut/gefahrgut_61.png}
\\\addlinespace\midrule\addlinespace
Klasse~6.2

Ansteckungs\-gefährliche Stoffe 
&
~\newline
\includegraphics[width=1.5cm]{./images/GEF/gefahrgut/gefahrgut_62.png}
\\\addlinespace\midrule\addlinespace
Klasse~7

Radioaktive Stoffe 
&
\includegraphics[width=1.5cm]{./images/GEF/gefahrgut/gefahrgut_7.png}
~
\includegraphics[width=1.5cm]{./images/GEF/gefahrgut/gefahrgut_7_fissile.png}
\\\hline
Klasse~8

Ätzende Stoffe 
&
\includegraphics[width=1.5cm]{./images/GEF/gefahrgut/gefahrgut_8.png}
\\\addlinespace\midrule\addlinespace
Klasse~9

Verschiedene gefährliche Stoffe und Gegenstände &


\includegraphics[width=1.5cm]{./images/GEF/gefahrgut/gefahrgut_9.png}
\\\addlinespace\midrule\addlinespace
\end{supertabular}

Quelle: ADR - Internationale Beförderung gefährlicher Güter auf der
Straße in der ab 1. Januar 2011 geltenden Fassung




\bigskip


 
\includegraphics[width=1.2917in,height=0.9583in]{./images/GEF/gefahrgut/warntafel.png}


  
\includegraphics[width=1.0937in,height=0.9583in]{./images/GEF/gefahrgut/erwaermter_zustand.png}

 
\includegraphics[width=1.0417in,height=1.0417in]{./images/GEF/gefahrgut/gefahrgut_umwelt.png}

\includegraphics[width=1.1043in,height=1.1043in]{./images/GEF/gefahrgut/abfall.jpg}


(Beispiel: Gefahrentafel für Benzin)\newline
Rückstrahlende orangefarbene\newline
Warntafel am Fahrzeug.\newline
40~cm x 30~cm. 

Kennzeichen für Stoffe,\newline
die in erwärmten Zustand\newline
befördert werden 

Kennzeichnung für umweltgefährdende Stoffe (Fisch + Baum) 


(rahmenloses Schild)\newline
Kennzeichnung für\newline
Abfallbeförderungen nach §~10 AbfVerbrG (Abfallverbringungsgesetz)


Obere Hälfte kennzeichnet die Gefahr
(33) - auch Kemlerzahl

Gefahrnummer 2 = Entweichen von Gas durch Druck oder chemische
Reaktion\newline
Gefahrnummer~3 = Entzündbarkeit von flüssigen Stoffen (Dämpfen) und
Gasen oder\newline
~~~~~~~~~~~~~~~~~~~~~~~~~~~~~~~~~~~ selbsterhitzungsfähiger flüssige
Stoffe (z.B. Benzin, Diesel)\newline
Gefahrnummer~4 = Entzündbarkeit von festen Stoffen oder
selbsterhitzungsfähiger fester Stoffen\newline
Gefahrnummer~5 = Oxidierende (brandfördernde) Wirkung\newline
Gefahrnummer~6 = Giftigkeit oder Ansteckungsgefahr\newline
Gefahrnummer~7 = Radioaktivität\newline
Gefahrnummer~8 = Ätzwirkung\newline
Gefahrnummer~9 = Gefahr einer spontanen heftigen Reaktion\newline
Gefahrnummer~0 = keine weitere Gefahr

Die Verdoppelung einer Ziffer weist grundsätzlich auf die Zunahme der
Gefahr hin. Einige Doppelziffern oder Dreifachziffern~ haben eine
besondere Bedeutung (z.B. 26 = giftiges Gas, 263 = giftiges Gas,
entzündbar, 33 = leicht entzündbarer flüssiger Stoff - Flammpunkt unter
23 °C, 70 = radioaktiver Stoff)

Ist der Ziffer ein {\textquotedbl}X{\textquotedbl} vorangestellt, so
reagiert das transportierte Produkt in gefährlicher Weise auf Wasser,
Nebel, Schnee (Löscharbeiten!).


Untere Hälfte bezeichnet den
transportierten Stoff~ (1203)\newline
(UN-Nummer), der von den Hilfsmannschaften aus einem Katalog zu erfahren
ist.\newline
Beispiele:\newline
0048 Sprengkörper\newline
0333 Feuerwerkskörper\newline
1005 Ammoniak\newline
1011 Butan\newline
1013 Kohlendioxid\newline
1017 Chlor\newline
1090 Aceton\newline
1133 Klebstoffe\newline
1203 Benzin\newline
1223 Kerosin\newline
1299 Terpentinöl\newline
1306 Holzschutzmittel, flüssig\newline
1779 Ameisensäure\newline
1789 Salzsäure\newline
1830 Schwefelsäure\newline
1962 Ethylen\newline
1972 Methan oder Erdgas (tiefgekühlt, flüssig)\newline
1977 Stickstoff\newline
1978 Propan\newline
2015 Wasserstoffperoxid\newline
2059 Nitrocellulose, entzündbar\newline
3295 Kohlenwasserstoff, flüssig\newline
3312 Gas, tiefgekühlt, flüssig














\section{Allgemeine Einsatzrichtlinie}
\label{topic:einsatzrichtlinie}
% TODO : Layout-Anpassung
\TODO{GABS}{yyyy-mm-dd}{Layout-Anpassung notwendig.}{Dieser Teil entspricht nicht dem Standard-Layout!}


\Layout{0}{Übersicht - GAS-Regel}{

\begin{itemize}
  \item \EE{G}efahr erkennen
  \item \EE{A}bsperrmaßnahmen durchführen
  \item \EE{S}pezialkräfte anfordern
\end{itemize}
}{
\begin{itemize}
  \item \EE{G}efahr erkennen
  \item \EE{A}bsperrmaßnahmen durchführen
  \item \EE{S}pezialkräfte anfordern
\end{itemize}
}{}{}{}



\Layout{57}{}{}{}{
\noindent\begin{tabular}{p{3cm}p{\linewidth - 7.6cm}p{3.1cm}}
\toprule
\noindent\TabHeading{{\color{ubuntured}G}efahr erkennen}
 &
\noindent Die Suche nach Warntafeln ist bei umgestürzten Fahrzeugen erschwert
 &
~

\includegraphics[width=3cm]{./images/noauth/gefahrgut-unfall.jpg}
\\
\noindent\TabHeading{{\color{ubuntured}A}bsperr\-maßnahmen
durch\-führen} &
Die Absperrung sollte großräumig  sein, genügender Platz für
die technischen Hilfeleistungen muss einkalkuliert werden. Beachte sich
ändernde Situationen  (Windrichtung, Niederschlag)
 &
~

\includegraphics[width=3cm]{./images/noauth/gefahrgut-absperrung.png}
\\
\noindent\TabHeading{{\color{ubuntured}S}pezial\-kräfte anfordern}
 &
Informationen angeben (Kemmler-Nummer, etc.)
 &
~

\includegraphics[width=3cm]{./images/gabriel-sebastian-aass/Atemschutz_IMG_2043_00800.jpg}
\\\bottomrule
\end{tabular}
}{}{}


\subsection{Gefahr erkennen}

{\TakeNotesLarge}


\subsection{Absperrmaßnahmen durchführen}
\label{topic:absperrmassnahmen}
% TODO : Layout-Anpassung
\TODO{GABS}{yyyy-mm-dd}{Layout-Anpassung notwendig.}{Dieser Teil entspricht nicht dem Standard-Layout!}

\TakeNotesLarge

\Layout{57}{}{}{}{
\begin{tabularx}{\linewidth}[H]{XXX}
\includegraphics[width=\linewidth]{./images/khd/gef-wind-rauch1.jpg}
\Captionof{figure}{\AuthNotesPicLegalOk{}{}Ein Auto brennt und alle schauen
fasziniert zu \ldots \SourcePicture{Gabriel}} 
&
\includegraphics[width=\linewidth]{./images/khd/gef-wind-rauch2.jpg}
\Captionof{figure}{\AuthNotesPicLegalOk{}{}\ldots doch dann kommt Wind auf \ldots  
\SourcePicture{Gabriel}} 
&
\includegraphics[width=\linewidth]{./images/khd/gef-wind-rauch.jpg}
\Captionof{figure}{\AuthNotesPicLegalOk{}{}Und oh Schreck! Die Absperrung
verliert ihren Zweck!  \SourcePicture{Gabriel}} 
\\\bottomrule
\end{tabularx}

\Fazit{Die Moral aus der Geschicht': Vertrau' dem lauen Klima nicht!}
}{Bilderserie \SlideHeading{Absperrung}.}{}

{\TakeNotesLarge}


\subsection{Spezialkräfte anfordern}
% TODO : Layout-Anpassung
\TODO{GABS}{yyyy-mm-dd}{Layout-Anpassung notwendig.}{Dieser Teil entspricht nicht dem Standard-Layout!}


\TakeNotesLarge
% TODO : Neufassung
\TODO{GABS}{yyyy-mm-dd}{Fehlender Text.}{}

\nocite{Cicero911CommissionReport}



